\section{总结与延伸}

Perceiver 的思想可以压缩成一句话：不要让巨大 raw input 在每层彼此两两注意，而是让少量 learned latents 读取输入、在内部深度计算，再用 task-specific output queries 解码需要的结构。这个重排让 attention 更容易扩展到像素、bytes、音视频与稠密向量场，也把架构设计问题转化为 latent capacity、position/modality features 和 query construction。

\begin{enumerate}
  \item \textbf{研究动机是系统可扩展性}：逐模态 handcraft 无法覆盖所有传感器与科学数据；
  \item \textbf{Transformer 的通用性来自弱连接假设}：non-local attention 与 position-as-feature 易迁移，但 raw self-attention 昂贵；
  \item \textbf{Latent bottleneck 改写复杂度}：输入处理从 $M^2$ 变为 $MN$，深层计算主要是 $N^2$；
  \item \textbf{线性 scaling 有条件}：latent 数固定、output decoder 可承受，且瓶颈没有丢失任务关键信息；
  \item \textbf{Perceiver 与 ViT 优化不同目标}：ViT 用 image patch prior 换效率，Perceiver 用较弱假设换跨域接口；
  \item \textbf{位置没有被删除}：Fourier/learned position 仍提供真实结构，只是不硬编码局部连接；
  \item \textbf{Output queries 定义任务}：位置、模态和 task embedding 告诉 decoder 应从共享 latents 读取什么；
  \item \textbf{Bytes 弱化 tokenizer 假设}：统一 vocabulary 的代价是更长序列和更多结构学习；
  \item \textbf{Optical flow 检验 structured output}：每像素 query 解码二维向量，synthetic-to-real transfer 展示弱偏置的潜力；
  \item \textbf{质性与定量必须并用}：EPE、PSNR、压缩率、吞吐和 failure slices 分别回答不同问题；
  \item \textbf{Generality 与 speed 是 Pareto tradeoff}：已知领域和小数据下，convnet/RAFT 等专用模型仍可能更优；
  \item \textbf{未完成方向很明确}：small-data learning、hierarchical cross-attention、联合多模态训练与更高效输出解码。
\end{enumerate}

\begin{importantbox}{课程作业：构建一个统一输入/输出实验}
选择两种结构差异明显的数据，例如图像 patches 与时间序列。固定 latent 数，使用共享 Perceiver processor，但分别设计 position/modality features 和 output queries。比较：专用 baseline、ViT/Transformer、Perceiver、不同 latent 数与去掉位置编码。报告质量、训练数据曲线、wall-clock、峰值内存、输入/输出长度 scaling 和 failure cases。
\end{importantbox}

\begin{warningbox}{最终自检：六个不能混为一谈的结论}
\textbf{弱偏置}不等于\textbf{无偏置}；\textbf{线性输入复杂度}不等于\textbf{线性端到端成本}；\textbf{input permutation robustness}不等于\textbf{空间结构无用}；\textbf{byte-level}不等于\textbf{零预处理}；\textbf{attention pattern}不等于\textbf{因果解释}；\textbf{统一架构接口}不等于\textbf{单一联合模型已训练完成}。
\end{warningbox}

\subsection{拓展阅读}

建议按“Transformer 基线—Perceiver 输入瓶颈—Perceiver IO 输出接口—领域对照”的顺序阅读：

{\small
\begin{itemize}[nosep,leftmargin=*]
  \item \href{https://arxiv.org/abs/1706.03762}{\textbf{Attention Is All You Need}}：QKV attention 与标准 Transformer 基线。
  \item \href{https://arxiv.org/abs/2103.03206}{\textbf{Perceiver}}：latent bottleneck、iterative attention、ImageNet 与多模态实验。
  \item \href{https://arxiv.org/abs/2107.14795}{\textbf{Perceiver IO}}：output queries、bytes、optical flow 与 structured outputs。
  \item \href{https://arxiv.org/abs/2010.11929}{\textbf{Vision Transformer}}：image patching 与视觉 Transformer 对照。
  \item \href{https://arxiv.org/abs/2005.12872}{\textbf{DETR}} 与 \href{https://arxiv.org/abs/2006.15055}{\textbf{Slot Attention}}：视觉 cross-attention / learned queries 的重要前身。
  \item \href{https://arxiv.org/abs/2006.10739}{\textbf{Fourier Features}}：高频坐标编码的理论与实验背景。
  \item \href{https://arxiv.org/abs/2103.06874}{\textbf{CANINE}} 与 \href{https://arxiv.org/abs/2105.13626}{\textbf{ByT5}}：tokenization-free / byte-level language 路线。
  \item \href{https://arxiv.org/abs/2003.12039}{\textbf{RAFT}} 与 \href{https://arxiv.org/abs/2104.14544}{\textbf{AutoFlow}}：optical-flow 架构偏置与 synthetic training 数据。
\end{itemize}
}

\end{document}
